\section{Zadanie 4}
Niech \(G\) będzie \(n\)-wierzchołkowym grafem i niech \(k\) będzie dodatnią liczbą całkowitą taką, że
\[k < \frac{n+2}{\Delta(G)+1}\]
Rozważa,y łańcuch Markowa, gdzie przestrzeń stanów \(\Omega_{k}(G)\) jest rodziną wszystkich \(k\)-elementowych zbiorów niezależnych w \(G\).
Krok ze stanu \(M_{t}\) do następnego stanu działa następująco.
Losujemy jednostajnie i niezależnie: \(v \in M_{t}, w \in V(G)\). Wtedy \(M_{t+1}=m(v, w, M_{t})\) czyli:
jeśli \(M_{t}\textbackslash\{v\}\cup\{w\}\in\Omega_{k}(G)\), to \(M_{t+1} = M_{t}\setminus\{v\}\cup\{w\}\); wpp. \(M_{t+1} = M_{t}\).
Udowodnij, że ten łańcuch jest nieprzywiedlny, nieokresowy i jego rozkład stacjonarny jest jednostajny.

\subsection*{Rozwiązanie}
Do nieprzywiedlność wystarczy pokazać że z \(a\) dasię dojść do dowolnego \(b\).
Pokażemy że z \(a\) dasię dojść do wierzchołka \(c\), który ma więcej wspólnych wierzchołków z \(b\).
Niech \(\{x\} \in b\) oraz \(\{x\} \notin a\) wtedy dopóki ilość krawędzi od \(x\) do wierzchołków z \(a\) większa od 1 to zamieniamy dowolny połączony z \(x\) na niepołączony\(^{*}\).
Jeśli \(\{x\}\) ma dokładnie 1 krawędź do \(a\) to zamieniamy po tej krawędzi, jeśli zero to zamieniamy \(x\) z dowolnym wierzchołkiem z \(a\) którego niema w \(b\).
To zawsze dasię zrobić jeśli \(y\) możemy zamienić na dowolny wierzchołek niepołączony z \(a\textbackslash{y}\cup{x}\),
ilość wszystkich wierzchołków w \(a\textbackslash{y}\cup{x}\) oraz ich sąsiadów będzie co najwyżej \((\Delta(G)+1)k - 2\), bo będzie co najmniej jedna krawędź w środku.
\[(\Delta(G)+1)k - 2 < (\Delta(G)+1)\frac{n+2}{\Delta(G)+1} - 2 = n + 2 - 2 = n\]
Co oznacza że istnieje co najmie jeden taki wierzchołek na który możemy zamienić \(y\).

Nieokresowość mamy z tego że da każdego z \(a\) mamy pętelkę do \(a\) kiedy \(v = w\).

Jeśli \(a, b \in \Omega_{k}(G)\) i istnieje krawędź z \(a\) do \(b\), to istnieje krawędź z \(b\) do \(a\).
Z tego że istnieje krawędź z \(a\) do \(b\), mamy że istnieją \(v \in a\) oraz \(w \in b\), \(v, w \in V(G)\), \(a\textbackslash\{v\}\cup\{w\} = b\).
Wtedy zamieniamy \(v\) z \(w\) miejscami i otrzymujemy krawędź z \(b\) do \(a\).
Także każda krawędź która nie jest pętlą ma to same prawdopodobieństwo i ono jest równe \(\frac{1}{k}\frac{1}{n}\).
Z dwóch powyższych faktów dostajemy że jeśli rozkład będzie jednostajny to on przejdzie sam na siębie, co oznacza że on będzie stacjonarnym.